\documentclass[10pt]{article}
\usepackage[T1]{fontenc}
\usepackage{lmodern}
\usepackage[a4paper,margin=23mm]{geometry}
\usepackage{amsmath}
\usepackage{booktabs}
\usepackage{array}
\usepackage{xcolor}
\usepackage{microtype}
\usepackage[hidelinks]{hyperref}

\definecolor{sourceblue}{RGB}{225,238,250}
\definecolor{errorred}{RGB}{252,232,230}
\definecolor{repairgreen}{RGB}{229,244,235}
\definecolor{softgray}{RGB}{245,246,247}
\setlength{\parindent}{0pt}
\setlength{\parskip}{0.55em}
\newcommand{\artifact}[1]{\texttt{#1}}

\title{When Executability Is Not Faithfulness}
\author{Explaining a Two-Stage Teacher Failure and Its Repair}
\date{}

\begin{document}
\maketitle

\begin{center}
\colorbox{softgray}{\parbox{0.90\linewidth}{\small
\textbf{Case.} \artifact{L2::bridge\_join\_0155\#L2a}.  A two-stage teacher produces an executable
program that returns 1,137.  Every online check passes, but the hidden benchmark answer is 63,118.
The successful repair does not change the named buyer or the aggregation operator.  It inserts the
missing target population between the derived CPV set and the final count.}}
\end{center}

\section{The question is about a derived population}

The benchmark asks:
\begin{quote}
How many procurement notices are listed under CPV codes that City of Doncaster Council has used
before?
\end{quote}

The named buyer is not the final population.  It defines a source population from which a set of CPV
codes must be derived.  Those codes then define a new target population over the entire procurement
snapshot.  The intended computation is
\[
\underbrace{\text{Doncaster records}}_{1{,}137}
\;\longrightarrow\;
\underbrace{\text{distinct CPV codes}}_{32}
\;\longrightarrow\;
\underbrace{\text{all records with those CPVs}}_{63{,}118}
\;\longrightarrow\;\text{count}.
\]

This distinction is easy to lose in natural language.  ``CPV codes that the council has used'' is a
set-valued condition inside the noun phrase modifying ``procurement notices.''  Counting the source
records or the codes answers a related but different question.

\section{Stage 1 identifies the bridge but chooses the wrong endpoint}

The Stage 1 model, \artifact{gpt-5.4-nano}, is asked to convert the question into a typed computation
program without answering it.  Its stored response is equivalent to
\begin{verbatim}
A = filter_records(buyer = City of Doncaster Council)
B = distinct(A, field = CPV code)
C = count(B)
return C
\end{verbatim}

The response gets two important facts right.  It grounds City of Doncaster Council as the buyer,
and it understands that the buyer's records produce a set of CPV codes.  The final type is wrong,
however.  The question requests a count of records, whereas \artifact{C} counts the intermediate
entity set.

Stage 1 also reports that ``used before'' lacks a temporal anchor.  In this benchmark family the
phrase denotes codes observed for the buyer in the fixed graph snapshot; it does not define a
per-record temporal comparison.  The model therefore introduces an ambiguity that is not part of
the executable benchmark contract.

\begin{center}
\colorbox{errorred}{\parbox{0.88\linewidth}{\small
\textbf{Stage 1 failure.}  The model discovers the source-to-CPV relation but treats the derived CPV
set as the answer population.  It also over-interprets a benchmark surface phrase as a missing
temporal constraint.}}
\end{center}

\section{Why Stage 2 is invoked}

The intent program can be compiled, but the committed routing code applies a semantic fast-path
check before accepting it.  For this row the check returns
\begin{center}
\artifact{count\_target\_is\_entity\_set\_for\_record\_question}.
\end{center}
This veto prevents the deterministic compiler from being the final planner and invokes the
configured lean Stage 2 planner, \artifact{grok-4-1-fast-non-reasoning}.  Stage 2 sees the original
question, the complete Stage 1 briefing, retrieved schema context, and the graph-planning rules.

The historical compact corpus preserves the Stage 1 briefing and the normalized rejected graph,
but not an explicit route tag or the raw provider request and response bytes.  The two-stage route
is therefore established by applying the committed routing code to the stored Stage 1 object.  The
same code emits the full reconstructed Stage 2 prompt included in the companion prompt bundle.

\section{Stage 2 remains executable but incomplete}

The normalized rejected graph contains a buyer-filtered record set and a derived entity variable:
\begin{verbatim}
A = record_set(buyer = City of Doncaster Council)
B = entity_set(depends_on = A, role = none)
return count(B)
\end{verbatim}

Stage 2 does not create the required target record set.  Its entity role is also left unspecified.
Execution consequently remains on the source side of the bridge and returns 1,137, the size of the
Doncaster record population, rather than the number of records sharing Doncaster's CPVs.

This is not a parser failure.  The graph has valid fields, a supported count operation, a grounded
buyer, a nonempty execution, and a numeric answer.  It is therefore a particularly informative
failure: the program is operationally valid while its population semantics are incomplete.

\begin{table}[h]
\centering
\small
\begin{tabular}{p{0.34\linewidth}p{0.16\linewidth}p{0.37\linewidth}}
\toprule
Contract & Verdict & What it establishes \\
\midrule
Schema and grounding & Pass & The fields and named buyer are executable \\
Operation support & Pass & Count is supported on the returned object \\
Execution and evidence & Pass & The program runs and returns nonempty evidence \\
Output consistency & Pass & The result has the expected numeric shape \\
Offline oracle agreement & Fail & 1,137 does not equal 63,118 \\
Constraint faithfulness & Exposed by audit & The target record population is absent \\
\bottomrule
\end{tabular}
\caption{Why an online success and an offline semantic failure can coexist.}
\end{table}

\section{The verifier does not have contradictory opinions}

The online and offline verdicts answer different questions.  The online runtime asks whether the
submitted action is grounded, executable, supported, and internally consistent.  The offline
evaluator additionally asks whether its result agrees with a benchmark oracle.

Let $p$ be the proposed graph program and $X(p)$ its deterministic execution.  The observed outcome
is
\[
f_{\mathrm{execute}}(p)=1,
\qquad
f_{\mathrm{release}}(X(p))=1,
\qquad
f_{\mathrm{oracle}}(X(p),y^*)=0.
\]
There is no contradiction.  Executability is a weaker claim than faithfulness to the full question.
The result illustrates the paper's execution--faithfulness gap:
\[
\text{successful execution}\;\not\Rightarrow\;\text{correct reasoning population}.
\]

\section{Repair changes the graph, not the answer string}

The offline controller uses the hidden oracle only to decide that a repair attempt may be collected.
The repair model is told that the submitted answer failed hidden validation and that the reference
answer is withheld.  It sees 1,137 as the submitted answer but does not see 63,118.

The accepted repair is
\begin{verbatim}
a1 = record_set(buyer = City of Doncaster Council)
b1 = entity_set(role = cpv, depends_on = a1)
c1 = record_set(depends_on = b1)
relation a1 -> b1, bind CPV
relation b1 -> c1, bind CPV
return count(c1)
\end{verbatim}

The critical change is \artifact{c1}.  It turns the derived codes into a filter over a new record
population.  Deterministic replay produces

\begin{center}
\setlength{\tabcolsep}{10pt}
\begin{tabular}{clrl}
\toprule
Variable & Kind & Size & Meaning \\
\midrule
\artifact{a1} & Record set & 1,137 & Doncaster's source records \\
\artifact{b1} & Entity set & 32 & Distinct CPV codes from \artifact{a1} \\
\artifact{c1} & Record set & 63,118 & All records whose CPV is in \artifact{b1} \\
\bottomrule
\end{tabular}
\end{center}

The repaired answer is therefore 63,118.  The same online checks pass again, and the result now also
agrees with the hidden oracle.

\begin{center}
\colorbox{repairgreen}{\parbox{0.88\linewidth}{\small
\textbf{Repair attribution.}  The improvement cannot be attributed to a new buyer grounding,
additional literals, answer copying, or a different aggregation operator.  It is attributable to
one inspectable structural change: inserting the CPV-bound target record population.}}
\end{center}

\section{What this phenomenon shows}

First, stronger generation alone is insufficient.  Both teacher stages can preserve many local
details while losing the global answer population.  Second, a verifier has bounded authority.  A
runtime verifier should not be described as proving semantic correctness when it has no independent
oracle for the intended population.  Third, an executable environment makes improvement
attributable.  The proposal, intermediate sets, feedback boundary, repair, and changed outcome can
all be inspected independently.

The claim supported by this case is consequently narrow but useful:
\begin{quote}
The environment does not create or certify improvement.  It makes reasoning actions, their
intermediate populations, and the feedback pathway observable enough to locate both a gain and the
boundary of verification.
\end{quote}

\section{Reproduction and evidence status}

All stored objects are joined by \artifact{L2::bridge\_join\_0155\#L2a}.

\begin{tabular}{p{0.49\linewidth}p{0.43\linewidth}}
\toprule
Artifact & Preserved evidence \\
\midrule
\path{data/qa/cicada_core_v4/train.jsonl} & Question, latent specification, oracle \\
\path{data/qa/teacher_full_v1/traces.jsonl} & Briefing, rejected graph, 1,137, verdict \\
\path{data/qa/teacher_full_v1/hard_negatives.jsonl} & Offline rejected example \\
\path{data/qa/teacher_full_v1/repair_sft.jsonl} & Bounded feedback and accepted repair \\
\path{data/qa/teacher_full_v1/dpo_pairs.jsonl} & Chosen and rejected plans \\
\path{doncaster_two_stage_teacher_prompt_bundle.txt} & Full readable prompt reconstruction \\
\path{doncaster_two_stage_teacher_prompt_bundle.json} & Structured prompts, schemas, and evidence labels \\
\bottomrule
\end{tabular}

The historical rows support the question, parsed Stage 1 output, normalized rejected graph, initial
answer, oracle mismatch, and accepted repair.  The complete Stage 1 and Stage 2 messages are emitted
again by the committed prompt builders.  Because raw historical provider messages were compacted,
they are labelled reconstructions rather than byte-identical API transcripts.  The
$1{,}137\rightarrow32\rightarrow63{,}118$ transition is recomputed on the versioned graph with the
accepted stored plan.

\end{document}
